\section{The Physics Validation: The \texorpdfstring{$E_8$}{E8} Substrate and the Geometric Cascade}
\label{sec:physics_validation}

As derived in \textit{Informational Energetics: A Universal Architecture of Persistence} \cite{meyer_informational_2026}, the vacuum is the unique resource-constrained computational substrate satisfying the six architectural pillars of persistence: an $E_8$ lattice projected to $D=4$ spacetime, with each localized node possessing strictly finite channel capacity ($\nu = 16$). This substrate contains four key invariants: the temporal resonance $\Delta=43$, the chiral capacity $\nu=16$, the interaction rank $\sigma=5$, and the topological boundary $\chi=2$, all derived from the $E_8 \rightarrow D_4 \oplus D_4$ projection.

While we derived the functional operator forms in \textit{Informational Energetics: Entropic Action} \cite{meyer_ieaction_2026}, the numerical coefficients remained undetermined. Standard QFT treats these constants as independent empirical inputs, each determined by measurement. In the IE framework, they emerge as coupled partition coefficients of a single finite capacity budget. Unlike $\alpha^{-1}$, which represents the unpartitioned baseline impedance derived directly from the substrate, the remaining parameters are not fundamental. They emerge strictly as the partition coefficients of the vacuum's finite channel capacity ($\nu = 16$): the unique allocation required for nested subsystems to satisfy global solvency.

If the spatial geometry, temporal updates, and internal force generators all attempt to saturate the node simultaneously, local complexity exceeds the 16 available channels and the causal structure becomes unresolvable. Because the vacuum persists, its structure must satisfy the three rules of Nested Persistence (\cref{sec:nested_persistence}). The $\nu = 16$ budget is therefore partitioned across orthogonal interaction sectors:

\textbf{Rule 1 (Capacity Partitioning)} allocates the $\nu = 16$ chiral capacity across orthogonal geometric sectors. This yields \textbf{System 4: The Gauge Partition}: the saturation ratio (Strong Coupling $\alpha_s$), the spacetime split (Weak Mixing Angle $\sin^2\theta_W$), and the irreducible geometric frustration of projecting 5-fold internal symmetry onto 4-fold spacetime (Jarlskog Invariant $J$). Because the temporal and color sectors are geometrically isolated by this partition, the Strong CP problem is structurally prohibited without requiring an axion.

\textbf{Rule 3 (Impedance Matching)} stabilizes the boundary between the high-frequency vacuum resonance and local topological defects. This yields \textbf{System 5: The Resonant Interface}, generating the electroweak mass scale: the Higgs VEV ($v$), Fermi constant ($G_F$), and the minimum resolvable mass (Electron Yukawa $y_e$), the floor below which vacuum fluctuations erase distinguishable states.

\textbf{Rules 1 \& 2 (Capacity Partitioning and Reversible Encoding)} govern the residual bandwidth, now attenuated across the geometric bulk. This yields \textbf{System 6: The Geometric Depth}, identified as Gravity. The attenuation depth derives the Planck scale ($M_P$); the irreducible bulk resolution floor derives the Vacuum Energy Density ($\rho_{vac}$), resolving both the Gravitational Hierarchy and the Vacuum Catastrophe.

\subsection{Note on Absolute Scale}

A discrete lattice inherently produces pure dimensionless ratios: $y_e = m_e/v$, $v/m_e \approx 4.8 \times 10^5$, $M_P/m_e \approx 2.4 \times 10^{22}$. These ratios are the framework's genuine zero-parameter predictions. To express them as absolute values in laboratory units (MeV, GeV), one scale must be fixed --- exactly as expressing the speed of light in m/s requires the definition of a meter. This anchoring is intrinsic to any scale-invariant framework and introduces no free parameter.

We adopt $m_e \approx 0.511\,\mathrm{MeV}$ as the \textbf{natural unit anchor}. This choice is not arbitrary: the electron mass is the \textbf{Landauer Limit} of the substrate, derived below in \cref{subsec:margin_electron} as the Margin output of System 5. It is the minimum resolvable topological defect --- any lighter charged state would dissolve into vacuum thermal noise. The electron is therefore the \emph{structural floor} that the partition rules necessarily produce; it is the natural unit in the same sense that Planck units are natural for quantum gravity.

All higher mass scales ($v$, $M_W$, $M_H$, $M_P$) follow as strict geometric multiples of this floor. The derivations do not use $m_e = 0.511\,\mathrm{MeV}$ as a numerical input; they derive the dimensionless ratios, and the single anchor converts them to laboratory units. This cascade spans 64 orders of magnitude from the gravitational coupling ($\alpha_G \sim 10^{-45}$) to the Planck mass ($M_P \sim 10^{19}$ GeV) with \textbf{zero tunable parameters}.